\section{Discussion and future work}
\label{sec:discussion}

\subsection{Paper-to-code is not paper-to-evidence}
A generated function can be syntactically correct and numerically stable while still using a misread table cell, an assumed unit or an unsupported convention. The cases show that the useful output of evidence automation is often not a match with a printed number. Three examples:
\begin{itemize}[leftmargin=1.5em]
  \item a demonstrated internal contradiction, as in P41 and P43;
  \item a named missing input, as in P44 and P45;
  \item a documented model-form difference, as in P42.
\end{itemize}
A system judged only by how often it reproduces printed numbers would be rewarded for normalising these away.

\subsection{Verification must itself be verifiable}
The most transferable lesson came from the framework's own verification. A check that recomputes in place destroys the distinction between reproducing evidence and replacing it. Only once verification was isolated and classified did cross-platform recomputation noise become visible. Only when it ran on further operating systems did platform-dependent provenance and an ill-conditioned independent check show up. Each case was then handled without touching frozen evidence. We expect the same pattern in any research-software evidence pipeline that reruns its own analyses.

\subsection{Related work and context}
Two strands of related work are relevant as context, not as validation evidence:
\begin{itemize}[leftmargin=1.5em]
  \item \textbf{Papers as agents.} Paper2Agent converts a research paper and its code base into an interactive agent: tools served through the Model Context Protocol that users can query in natural language \cite{Miao2026Paper2Agent}. \engiproof{} exposes its callable methods through the same protocol, but it adds an evidence envelope and a qualification boundary to every result.
  \item \textbf{Chains of evidence.} ScientistOne introduces \emph{Chain-of-Evidence}: every claim produced by a research system must be traceable, through a recorded chain of supporting claims, to a grounding source \cite{Meng2026ScientistOne}. ScientistTwo extends this line to a fully autonomous multi-agent research pipeline \cite{Nam2026ScientistTwo}.
\end{itemize}
They are related to \engiproof{} at the level of architecture. The contribution here differs in three ways:
\begin{itemize}[leftmargin=1.5em]
  \item the engineering-specific separation of published, independent and solver evidence;
  \item preservation of source inconsistencies, with a human decision boundary;
  \item non-mutating verification with an explicit qualification boundary.
\end{itemize}
\engiproof{} shares the traceability principle with this line of work but deliberately keeps a human decision boundary. Neither strand provides engineering validation of the P40--P45 evidence. \todo{Add reproducibility literature in computational science and engineering.}

\subsection{Future work}
\begin{itemize}[leftmargin=1.5em]
  \item \textbf{Solver adapters.} External solvers generate \eclass{SOLVER\_NEW} evidence on existing cases. The agent orchestrates the solver runs, and \engiproof{} adjudicates their results.
  \item \textbf{Extraction benchmark.} A labelled extraction benchmark would make accuracy measurable.
  \item \textbf{External replication.} Replication of selected studies by independent teams.
  \item \textbf{Experimental evidence.} Recording experimental reference data as a separate population where the sources contain it.
  \item \textbf{Environment metadata.} Frozen artifacts that declare the environment that produced them.
\end{itemize}
